\section{Introduction}
\label{sec:introduction}

Social media have been recently sentenced 
for their endangerement of the mental health of young people.
In the Netherlands, this has led to an imposition on Meta 
to reintroduce a reverse chronological feed~\cite{reuters-dutch-meta}.
This feed is also widely used by alternative social media,
like Mastodon, PeerTube and BlueSky.
However, even though it is extremely interpretable,
the reverse chronology recommendation algorithm fails 
to provide a satisfactory user experience.
Additionally, it favors hyperactive publishers,
and incentivizes fast and massive content creation.
More generally, it does not fully address the demands for influence traceability.

In this paper, we propose \Veche{}\footnote{
    Veche was a popular assembly among some Slavic peoples during the Middle Ages. \\
    The name is not set in stone yet, but I would strongly want it to promote democracy.
    Other candidates include T\'uath, Ecclesia or Landsgemeinde.
    Ideally, I'd love something close in spirit to the algorithm
    (i.e. selecting who you listen and gathering them to vote on what you should know).
    Feel free to check \url{https://en.wikipedia.org/wiki/History_of_democracy}
    or other sources to come up with other proposals.
},
an alternative, modular and flexible following-based recommendation algorithm
grounded in social choice theory principles.
\Veche{} implements the following key principle.

\begin{center}
    \emph{A follow is the allocation of a unit voting right of influence.}
\end{center}

Etymologically, a \Veche{} was a Slavic popular assembly,
the last of which was held in 1510.
We propose to select a user's recommendations by emulating such an assembly.
More specifically, by choosing to follow another user,
Alice is metaphorically inviting the followed user to act as a \emph{councillor}
in Alice's personal \Veche{}.
When Alice wants to be suggested entities (e.g. posts, videos or other kinds of options),
she asks here councillors to submit entity recommendations.
Each councillor may even specify the extent to which they want a given entity to be shown to Alice.
Such weighted entity recommendations thereby form the councillor's \emph{ballot}.
Alice's personal \Veche{} will then collect her councillors' ballots,
and construct out of them a personalized feed of entities,
which will be shared with Alice.

% \Veche{} considers that, 
% when Alice follows some user $u$,
% the user $u$ is given a unit voting right in a social choice mechanism
% that will then (stochastically) construct Alice's recommendation feed.
% Put differently, the user $u$ then gets to submit a ballot $b_u$.
% The user $u$ thereby gets to be recommendation \emph{councillors} 
% to the recommendation receiver Alice.
% A voting system will then aggregate all councillors' ballots,
% and return a feed based on these ballots,
% with the guarantee that each councillor $u$'s influence is bounded by their unit voting right,
% no matter what ballot $b_u$ is submitted.

Given that ballots must be be submitted upon any of Alice's query for a feed,
and given that she may be expecting to receive a feed of recommendations within a fraction of a second,
the construction of the councillors' ballots will have to be automated, 
based on councillors' past activities (posts, reposts, reports...).
This automated ballot-construction mechanism may called an \emph{algorithmic representative},
following the principles put forth 
by~\cite{DBLP:journals/pacmhci/LeeKKKYCSNLPP19,DBLP:conf/aaai/NoothigattuGADR18,DBLP:journals/ai/FreedmanBSDC20}.
In principle, a councillor could construct any algorithm to submit a ballot in their name.
However, in practice, most councelors will require a default construction of an algorithmic representative.
The paper also proposes such a default algorithmic representative.

More concretely, our contributions are the following.

\begin{itemize}
    \item We introduce \Veche{}, a highly flexible social choice mechanism,
    through which councillors can vote for the content of Alice's feed,
    and we provide mathematical guarantees on the maximal impact of a councillor.
    \item We show how to leverage \Veche{}'s flexibility and algorithmic representatives
    to propose a general and practical deployment that responds to user needs and constraints,
    such as minimal user experience modifications,
    simple time decay management with desirable mathematical properties,
    tunable councillors' volumes,
    including preferred semantic biases,
    encouraging reposts and reports over posts,
    and feed diversification using Tillé's pivot algorithm~\CITE{}.
    \item We provide mathematical proofs that \Veche{} 
    and some of its proposed instantiation details are well-behaved,
    thereby providing interpretability and trustworthiness.
    \item We provide a free and open-source implementation of \Veche{},
    as part of the Solidago library~\cite{hoang2022solidago}.
    We also deploy \Veche{} on Tournesol~\cite{hoang2025tournesol} and BlueSky~\cite{kleppmann2024bluesky}.
    We also provide the materials to effortlessly 
    adapt and redeploy our algorithms through the AT protocol~\cite{kleppmann2024bluesky}.
    % \item We evaluate our algorithms through user surveys,
    % thereby directly assessing their concrete impacts.
\end{itemize}

Overall, we believe that the desirable properties of our algorithm,
especially in terms of the traceability of influence,
make a compelling case to make it, perhaps through legal enforcement, 
the default recommendation algorithm for all social media,
instead of the reverse chronological feed.

The rest of the paper is organized as follows.
First, Section~\ref{sec:context} provides some context and a review of related literature.
Next, for mostly pedagogical reasons,
Section~\ref{sec:model} opens with a simple instantiation of \Veche{}, 
and provides proofs of its main guarantees.
Section~\ref{sec:generalizations} then shows generalization of the first instantiation,
to highlight the flexibility of \Veche{},
and to show how to respond to various needs in practical deployments.
Section~\ref{sec:evaluation} evaluates \Veche{} and other recommendation algorithms
along Tournesol's four digital democracy pillars.
Section~\ref{sec:experiments} reports our empirical analysis.
Section~\ref{sec:conclusion} concludes.
